\documentclass[11pt,a4paper]{article}
\usepackage{fontspec,xeCJK,geometry,needspace,longtable,array,amsmath,fancyhdr}
\usepackage[unicode,hidelinks,breaklinks]{hyperref}
\setmainfont{Microsoft YaHei}[BoldFont=Microsoft YaHei Bold]
\setCJKmainfont{Microsoft YaHei}[BoldFont=Microsoft YaHei Bold]
\newfontfamily\titlelatin{SimHei}[AutoFakeBold=2]
\newCJKfontfamily\titlecjk{SimHei}[AutoFakeBold=2]
\geometry{margin=22mm}\pagestyle{fancy}\fancyhf{}\fancyfoot[C]{SRP 新实验与分析候选\quad\thepage}\renewcommand{\headrulewidth}{0pt}
\setlength{\parindent}{0pt}\setlength{\parskip}{7pt}\setlength{\tabcolsep}{4pt}\linespread{1.22}\emergencystretch=3em\raggedbottom
\newcommand{\reviewtitle}[1]{\begin{center}{\titlelatin\titlecjk\bfseries\fontsize{21}{28}\selectfont #1}\end{center}}
\newcommand{\reviewsection}[1]{\par\Needspace{6\baselineskip}\begin{center}{\titlelatin\titlecjk\bfseries\fontsize{14}{20}\selectfont #1}\end{center}}
\newcommand{\reviewsubsection}[1]{\par\Needspace{4\baselineskip}\begin{center}{\titlelatin\titlecjk\bfseries\fontsize{12}{18}\selectfont #1}\end{center}}
\newcolumntype{L}[1]{>{\raggedright\arraybackslash}p{#1}}
\begin{document}
\reviewtitle{SRP实验与分析操作候选}
设计候选 ODES-20261008；与05控制规格、06配对故事板共同使用。研究只比较固定A与数据驱动原生受限适应B整包。本轮没有招募、试验、参数冻结或正式预注册；现行协议保持原范围。

设计审阅结论：2026年10月9日，用户明确“批准设计通过”。本轮设计方案及材料获批，作为后续前置工作依据。测量缺口M01/M02仍按既定程序处理；本批准不代表中文工作稿已验证、参数及正式样本量已冻结，或前置活动与正式研究已获实施准入。

\reviewsection{一、研究问题及能够回答的结论}
\begin{longtable}{L{16mm}L{40mm}L{63mm}L{35mm}}
问题 & 构念 & 测量/分析候选 & 结论边界\\\hline
RQ1 & 即时负性情绪体验 & PANAS-NA后测，前测校正；参与者级两臂分配效应及95\%区间 & 整包相对效益，非长期或单机制效益\\
RQ2 & 共同目标达成、实际引导跟随、努力和控制成本 & 独立X对I与有效G分报，覆盖/缓冲/偏移及负担 & 跟G好不代证达I；无事前界不称非劣\\
RQ3 & 何时支持/干扰及表示边界 & 预定状态案例、事件轨迹和体验后访谈；正反例对照 & 解释性证据；触发者与未触发者差异非因果\\
\end{longtable}

候选贡献是受限支持与目标漂移的可信权衡及有限设计知识，不宣称首次数据驱动、首次原生或新控制算法。最近工作差异仍由科学附录与W-01继续复核；独立迁移未执行时不能称原则已迁移。

\reviewsection{二、条件、教学与共同时间}
\reviewsubsection{2.1 保持与处理差异}
两组共同I、天气底图/镜头、非表现环境音、设备测量、操作员关注、帮助/停止规则、教学预算与模块暴露预算。A数据不能影响输出；B允许真实响应、累计练习响应和05中的受限时序。完成周期数和视觉轨迹可以不同，不给B补时补次数。

两组教学使用中性称呼，不称B高级。工作候选为180秒教学与800秒核心体验，保留为可行性剂量；四天气若保留，则每模块200秒是等预算候选而非最佳剂量。正式前确认四结构用途、舒适性及可辨性；不因已有资产而保留不可解释结构。

公共教学包括：自愿跟随、停留不必强求、漏步骤不追补、可以自然呼吸或停止。A另说明圈是预设引导、场景变化不是个人表现；B另说明目标、真实响应及信号不可用。各自内容准确、预算相同，不以教到相同正确率作为入组条件。检查题不与教学逐题复用。

\reviewsubsection{2.2 流程及提前停止}
中性准备/设备校准→条件揭示前PANAS前测→隐藏分配揭示→条件教学→共同预算体验→PANAS后测→努力/舒适→条件适配理解题、SCCI及访谈。情绪后测始终先于理解考试。

主要后测候选在计划体验结束时点进行。提前停止者可自愿自然休息至该时点；绝不要求为测量继续练习或留场。如立即离开，离开前的自愿情绪记录只作附加时间点，不替主要后测，该主要值按缺失处理。实际测量延迟单独记录。若该安排在前置阶段不可接受，正式前共同修订测量时间与估计目标。

暂停不计有效暴露，单报墙钟；自然缓冲和中性引导占暴露T。结束前不足以完成下一完整周期时，用剩余预算自然呼吸收尾，不赶做半个周期。退出者、未触发适应者及降级者保留其原分配，不按结果重新归组。

\reviewsection{三、测量蓝图与独立真值}
\reviewsubsection{3.1 主要量表及其他体验量}
PANAS-NA保留为主要候选。2026年10月9日起，设计采用已下载PANAS-X原始手册中的一般NA与PA分量表，各10题，共20题；不增加到60题，也不增加新的主要终点。英文原件可按非商业用途使用；正式中文题本仍须完成词义、来源与使用条件核对。

NA和PA各自按10项1--5相加，范围10--50，不反向计分，不合成正负情绪总分。任一分量表缺题则该分量表总分缺失，这是本项目工作规则，不冒充作者缺题规定。此刻指令由手册印刷页3支持；校园样本的理解、地板分布与短时变化适用性仍需前置验证。

原始手册为Watson与Clark，版权1994，1999年8月更新；已从爱荷华大学机构库取得完整28页PDF（正文24页），\href{https://doi.org/10.17077/48vt-m4t2}{DOI 10.17077/48vt-m4t2}。印刷页1允许按研究需要选用分量表，页2表2给出NA/PA题项归属，页3包含此刻指令。使用本手册不等于取得某一已验证中文译本。

努力与呼吸舒适分别用清楚标注的0--10单题候选，公开题干和两端含义；它们是自编过程题，不称已验证量表，不合并为总分。求助、停止、再次使用意愿及准备成本分别记录。PANAS-PA可作预先指定次要结果；SCCI只作融入操纵检查，不能证明数据实际驱动或情绪收益。

\reviewsubsection{3.2 独立X的构建}
原始胸带波形、质量/运动信息和采集时间是输入；独立判定方法冻结后以同一规则应用两组。I评分器不读取G、在线成功标志、场景状态或分组。将代表性正常/异常片段交给不知条件的独立标注者，报告一致性、错误、延迟与不可辨比例；不是用当前控制器输出循环验证控制器。

胸带振幅不能自动解释为绝对吸气量，低变化不能直接证明有意停留，第二吸气也不能从补潮动画推断。可信波形仍可能步骤不可辨，分别记录技术可用与步骤可判定。Fade补吸识别门未通过时，不开展声称完整三步适应的正式比较；先简化/修订范围并重新冻结。

\reviewsubsection{3.3 对I、对G和未知分开计分}
按独立X识别实际完整周期的起止和步骤顺序，再与对应天气的I模板比较。允许周期起点整体平移，禁止任意伸缩或时间扭曲；步骤次序缺口不能通过匹配消除。周期结束跨出模块边界时不重复归属：属于起点所在模块，预算内已发生事实保留，边界外不补入有效达成；末尾未完成单列，不伪作技术缺失。

先报告实际周期数、缺步骤、各步骤时长偏差、周期时长偏差及可判定覆盖。容差待独立误差与有意义偏离共同校准。完全满足冻结I结构/时长要求的周期，在T内的时间贡献记为C；周期互不重叠，同一时段不重复累计。缓冲中的自然呼吸若独立满足I也可贡献C，不能按B模式删掉或额外奖励。

技术可用时间与结果可判定时间分开：O\textsubscript{sig}表示原始信号可用时间；O表示能确定是否符合I的完整结果时段。T按同一独立方法划为已证实合格C、已证实不合格K及结果未知U，C+K+U=T，O=C+K。若短暂断流使整个周期不能判断，U覆盖整个不确定周期及必要边界区间，不能只用断流秒数；未知跨越的可见片段也不能草率记K。

候选可证实达成下界为C/T，上界为(C+U)/T；另报O\textsubscript{sig}/T、O/T及C/O（O=0时不可估计）。该界只针对这一操作定义，不把未知当失败或当成功，不证明未知段真的存在完整合格周期。已知未完成不纳入U。使用共同T，不能按B减少的周期机会改分母。

例：T=100、C=20、K=40、U=40，达成界为0.20--0.60，即使原始断流仅10秒，也不能把上界写成0.30。这个例子只验证计分逻辑，不是本项目观测结果。边界扩展与分区规则须在独立标注材料上验证后冻结。

候选分区程序：先由独立X算法给出完整实际周期边界及未确定事件，不读G/条件；每个不确定区间扩展到前后可靠的实际周期边界，取并集并裁到模块预算内。没有可靠边界时相应模块全段列U。跨越该并集的周期不能拆出可见部分记K；完全在其外的完整周期按共同I分类为C或K，其余已知未完成时段为K。核查三类互斥且覆盖T，不重复计时。离线核查可以使用后续原始波形帮助确定边界，但这些后见事实不能反流用于控制决策。

对G另用有效呈现时间线测步骤/同步误差，称实际提示跟随。G参考可用于该过程评分，不能倒流改变X标注或I评分。多天气结果以相同预定模块权重形成参与者级结果；不把每周期或每题当独立受试者。

\reviewsubsection{3.4 控制成本及信息读取}
控制成本独立包括：吸/呼时长相对I偏移的时间积分、最大尺度、偏移持续时间、自然缓冲占T、未回归比例与结束时偏移。停留时长不扩大；意外扩大是违反设计的事件。另报质量降级/中断、实际完整周期量及墙钟准备成本，避免用一个任意综合分遮盖代价。

理解蓝图：两组共同当前/下一目标及停止规则题，经等值验证后有限比较；A专属题核查是否误以为预设画面代表自身成功；B专属题核查目标/实际、累计与不可用含义。状态片段同源配对、真值预标，答案缺失与“不知道”分记。旧八题或新题数均不是正式主量表；理解只解释体验后的信息读取，不代证在线理解、情绪或长期学习。

\reviewsection{四、分析、缺失与解释规则}
\reviewsubsection{4.1 主要估计与模型}
估计B减A的分配效应，负值表示B的NA后测更低。候选模型
\[\mathrm{NA}_{post}=\beta_0+\beta_1 B+\beta_2(\mathrm{NA}_{pre}-\overline{\mathrm{NA}_{pre}})+\mathrm{order\ strata}+\varepsilon.\]
参与者为单位，预指定天气顺序效应，报告差异、双侧95\%区间及HC3方差。前测校正方法依据\href{https://www.bmj.com/content/323/7321/1123}{Vickers与Altman的原始统计方法文}；本轮出版社PDF请求被拒，已有题录/方法摘要核查不等于阅读全文。新模型还需针对分布与缺失验证。

不校正处理后的努力、理解、适应次数、剂量、美感或可信度。没有无项目组，不证明相对自然休息的绝对效益；没有机制分解条件，不分离原生、反馈和适应的贡献。

\reviewsubsection{4.2 缺失及多结果}
主分析按原分配；主要后测缺失以预指定多重插补候选处理，包含条件、前测、顺序及预先记录的预测变量；附完全案例和按两臂分别偏移缺失结果的不利方向敏感性。插补模型、迭代/重复数及偏移域需A-03验证并正式冻结，不能以前测代填冒充保守。若缺失导致不可辨识，应报告而不替换终点。

原始X未知不用G插补，报告覆盖与上下界。周期/天气重复观察要聚合到参与者，或使用显式层次模型；不扩大N。主要结果只保留一项，其他结果列次要/探索并如实说明多重性；不事后挑显著项组成“成功”。没有容差界、精度和预先分析集，不宣称功能非劣或替代性。

\reviewsubsection{4.3 五种结果均有出口}
情绪有相对收益且I/负担可接受：支持该场景中的整包相对价值。情绪有收益但I明显漂移：报告支持与目标放宽的权衡。无明确情绪差异：按区间与精度解释，不称等效；功能改善只支持对应功能结论。B更差：保留A并定位成本。信号不可辨/迁移失败：收缩闭环或设计原则范围，不能靠动画或重新命名维持贡献。

\reviewsection{五、人数、顺序与实际可行性}
\reviewsubsection{5.1 前置阶段按唯一人员记账}
\begin{longtable}{L{34mm}L{25mm}L{72mm}L{23mm}}
阶段 & 工作候选 & 要解决什么 & 计数\\\hline
专家/独立重建 & 8位评审候选，其中2位尝试独立重建 & 构念、规则、步骤用途与可重建；不是用户效益 & 8人，不再加2\\
需求/认知/舒适 & 12学生，条件各6；8人先做需求访谈 & 是否需要支持、两组教学、状态理解及负担 & 12，不另加8\\
真实技术与完整流程 & 12--20名新学生，默认规划20 & 真值、延迟、动作域、失效恢复与量表分布 & 与上阶段不同\\
正式整包研究 & 资源剩余且满足精度要求的N & RQ1/RQ2确认或估计研究 & 前置成员不转正式样本\\
独立新实例迁移 & 后续可选，另计样本/资源 & 有限规则是否能应用 & 不占本轮核心预算\\
\end{longtable}

若资源为100名学生，12认知+20技术后正式剩68；若同一100人总额还包含8位专家，正式剩60。若技术采用12人，分别剩76与68。希望正式100学生，默认前置需另32学生与8评审。100+的实际含义须团队明确，不重复计算同一人，也不靠减少必要前置来凑正式N。阶段数是工作量候选，不是效度普适阈值；技术阶段不足时增加前置并重新评估正式资源。

\reviewsubsection{5.2 精度量级，不倒填效应}
两臂等分、残差标准差为$\sigma$、双侧0.05的正态近似：
\[h_{95}/\sigma\approx3.92/\sqrt N,\qquad d_{80}\approx5.603/\sqrt N,\qquad d_{90}\approx6.483/\sqrt N.\]
\begin{center}\begin{tabular}{rrrr}\hline
正式总N & 95\%区间半宽/$\sigma$ & 80\%功效差/$\sigma$ & 90\%功效差/$\sigma$\\\hline
60 & 0.506 & 0.723 & 0.837\\
68 & 0.475 & 0.679 & 0.786\\
76 & 0.450 & 0.643 & 0.744\\
100 & 0.392 & 0.560 & 0.648\\\hline
\end{tabular}\end{center}
N指可用主要结果的独立参与者总量，不是题数、周期数或排列数。上述不是预期效应或正式功效，未包含小样本、协变量估计、不平衡、地板与缺失代价。先确定产品有意义的NA原始分差，再按预试残差/缺失区间模拟实际模型。若资源不能辨别该差，增加资源或明确定位为估计性研究，不把关注差改大来宣称充足。

\reviewsubsection{5.3 随机化候选取代“必须满48格”}
正式保留四天气时，候选四条Williams顺序为：Storm--Heat--Fade--Snow；Heat--Snow--Storm--Fade；Snow--Fade--Heat--Storm；Fade--Storm--Snow--Heat。每种天气各占每个位置一次，12种有向相邻对各出现一次。这是位置/一阶顺序平衡，不保证所有更高阶效应消失。

两条件乘四顺序形成8个分配格；候选隐藏联合分配块为8或16，比例1:1，独立生成并隐藏后续序列。末尾不完整块按已生成序列执行，报告实际不平衡，不因结果或“凑格”选择追加者。原48块依赖24种全排列，不是统计样本依据；12人批次也不是随机化块。X-01另实现并验证，当前旧算法不因此改变。

\reviewsubsection{5.4 五人、半个月的用途}
五职责分别承接科学/测量、设备真值、控制与日志、Unity呈现、用户流程/统计整合，独立复核交叉安排。半个月只评估后续实施可行性：用实测准备、校准、会话、重置及异常处理人时估算吞吐，不承诺所有前置和正式样本在期限完成。单会话工时用准备+校准+前测+教学+体验+后测/访谈+重置+暂停计算，不能只按800秒乘N。

\newpage\reviewsection{六、前置放行、查新与任务承接}
\reviewsubsection{6.1 需要完成的具体证据}
需求与多步必要性经真实学生和评审核查；A不误导、B关键状态可理解；Fade补吸及停留具有独立可辨证据；规则域具有舒适性与技术依据；A输入隔离、B合法动作、失效/恢复与实际呈现可重建；量表版本/地板、最小关注差与N有依据。每项先写接受条件，再执行验证；预试不以小样本组间显著性放行效益。

可行性与正式效益分开，依据\href{https://doi.org/10.1136/bmj.i5239}{Eldridge等先导与可行性报告方法}。本轮出版社全文访问受限；已有项目科学附录与可核查题录支撑方法定位，不据此给本项目人数普适门槛。量表版本、控制数值、容差、最小关注差及正式N仍未冻结。

\reviewsubsection{6.2 设计交付不替代真实签收}
F-02/Q-01/Q-02确认构念与需求；U12-02/U12-03完成条件适配材料和教学；Q-03/E-03执行技术校准；U12-04/A-03形成SAP与真实精度论证；X-01实现新随机化候选；U12-11/G-03冻结正式版本后才进入B-02。W-01继续最近工作核对，A-06/W-02按实际证据关闭贡献。

主终点、测量、随机化与阶段变化都是下一版协议候选。原状态、签收及领取快照保留，已领取包通过显式补充与影响记录消费本稿；本轮不修改运行或研究权限。

% NORMALIZATION-20261008
\newpage\reviewsection{七、可使用教学与测量材料}
\reviewsubsection{7.1 当前采用的原件、评分键与中文缺口}
黄丽、杨廷忠、季忠民，2003，17(1):54–56，DOI 10.3321/\allowbreak{}j.issn:1000-\allowbreak{}6729.2003.01.018的身份与摘要已在期刊站核准。侧面取得CHEUNG, Sin Yan Ziv的2018香港教育大学博士论文：方法印刷页41引用该译本，英文附录印刷页127/PDF138，中文印刷页133—134/PDF144—145含20题及五级评分。原件已下载并逐页目视核对。\href{https://www.lib.eduhk.hk/pure-data/pub/202201871.pdf}{机构库原件与附录}

该附录中英题序对应的NA候选索引为2、4、6、7、8、11、13、15、18、20；PA为其补集。无反向题；NA10题1至5合计10至50仅是对应20题结构的计分候选，不能套给18题修订。附录第13项ashamed对应“害羞”，需要核查语义及与原始译本的一致性，不自行改词后称为原版。

中英文附录均明确询问过去1—2周，而本项目需要此刻感受。现已取得直接使用先例：\href{https://doi.org/10.1111/aphw.12133}{Hu等2018研究，PDF第7页}采用完整20题中文版、此刻指令、五级评分及10题分量表求和，样本中PA/NA内部一致性分别为0.87/0.88；\href{https://doi.org/10.3724/SP.J.1041.2021.00244}{刘宇平等2021研究，PDF第7页}明确使用中文版NA10题及此时此刻指令。这支持状态测量的设计依据，但研究人群、任务和刺激不同，不能外推到普通校园间隙的地板分布或短时变化能力。

M01现缩为唯一中文题本、词义与验证来源逐项对应；第十节已提供内部审阅工作译稿，selected\_form与正式item\_texts仍为空。此刻作为有文献依据的工作指令保留；最终题本若有改词或指令适配，明确记录。在认知/前置验证中检查理解、分布、缺失和短时变化适用性，这是待取得的真实证据，不再将“状态测量没有使用先例”列为设计缺口。

馆藏页水印限制进一步复制/出版；本包不转载完整题文。另一\href{https://vtechworks.lib.vt.edu/bitstreams/6844cae0-de23-4429-afc9-12f4c01bc260/download}{机构库英文此刻版，附录O印刷页62/PDF70}中的特定转载许可不能转移给SRP。但这也不能推导所有研究施测都必须逐次申请：\href{https://www.apa.org/pubs/databases/psyctests/faq}{APA官方FAQ}说明，具体记录标为可用于研究/教学时允许相应非商业用途，公开题文及其他用途另行核对。\href{https://doi.org/10.1037/t03592-000}{原20题PANAS数据库记录}的身份与描述已经直接核准，具体许可字段尚未取得；中文译本也未核准。M02分别记录研究施测、改编与完整题文公开的适用条件，不把机构库水印当成量表本身的完整权利结论。F-02/U12-04承接M01，F-02/G-05承接M02；未核准前不形成正式主要测量材料。



补获第二份完整题本侧证：\href{https://pdf.hanspub.org/ap20220600000_21602996.pdf}{陈勇达、李文杰2022出版社PDF}第11页（印刷页2037）列出20题、五级选项，第13项也为“害羞”；附录未写明时间窗。该文引用黄丽等2003，但不能据此证明逐字等同原译本。文章第2页标示CC BY 4.0，尚不足以独立核准第三方量表全部用途。已在线读取，未取得本地PDF副本。

\href{https://fieldinglibrary.wordpress.com/wp-content/uploads/2014/09/psyctests-on-psycnet.pdf}{APA编写的PsycTESTS使用指南第10页}明确展示Chen, Yiwei、Peng, Yisheng、Fang, Ping（2016）的中文PANAS记录\href{https://doi.org/10.1037/t58959-000}{10.1037/t58959-000}：20题、过去一个月、独立翻译回译。它与优先核查的2003译本不是同一版本，不自动替换。图书馆检索应分别提供原版t03592-000及中文t58959-000的题本、来源与Permissions字段，不以“搜到PANAS”作为完成。

后续取得路径已经具体化：优先向Hu等2018论文通讯作者索取实际施测中文20题、此刻指令、计分及版本来源，并询问适用使用说明；联系人由论文首页核准为\texttt{zhangwc@psych.ac.cn}。索取函及图书馆检索单保存在\texttt{agent/work/2026-10-08\_科学来源补获与量表侧证.md}第六节。函件尚未发送；取得完整回复后由F-02与U12-04逐项核对，G-05核对使用范围，再决定唯一题本。若只有题本而无条件说明，M01与M02分别判定，不捆绑宣布关闭。

2026年10月9日补核：\href{https://www.phenxtoolkit.org/protocols/view/180501}{PhenX成人情绪协议180501}明确允许PANAS-X非商业使用，并列出中文下载入口。实际下载的\texttt{PX180501\_Chinese.pdf}仅1页（印刷页205），为中文协议说明，未提供60个中文题项、完整指导语或计分键。已保存并目视核查，不能将中文入口描述为已取得完整中文题本，也不能将扩展版许可转用于既有20题译本。

另核准\href{https://zxws.cbpt.cnki.net/portal/journal/portal/client/paper/3a515b38b5ea9848b85c2a7b8cb00a99}{郭明珠、甘怡群2010，24(7):524--528}的期刊摘要：660名大学生研究后的修订版为40题、9因素。它不是PhenX说明的原60题版本；本轮未取得其完整题本和译本使用条件。现采用PANAS-X原始手册的一般NA10与PA10作为20题设计来源，主要构念和前后测分析不变；该40题修订版不替代当前来源。完整中文施测题本仍待核准。

本次已落实的使用决定：按\href{https://www.phenxtoolkit.org/protocols/view/180501}{PhenX 180501非商业使用说明}，在本项目非商业研究设计中采用英文PANAS-X一般分量表。保留作者及版权标识：Copyright 1994 David Watson and Lee Anna Clark；手册updated 8/99。原件保存于项目本地来源目录的\texttt{PANAS-X-manual.pdf}。下表为印刷页2表2的评分归属，行号仅用于审阅，不规定正式问卷题序，也不是中文施测表。

\begin{longtable}{L{15mm}L{60mm}L{60mm}}
行 & NA（主要候选） & PA（次要候选）\\\hline
1 & afraid & active\\
2 & scared & alert\\
3 & nervous & attentive\\
4 & jittery & determined\\
5 & irritable & enthusiastic\\
6 & hostile & excited\\
7 & guilty & inspired\\
8 & ashamed & interested\\
9 & upset & proud\\
10 & distressed & strong\\
\end{longtable}

原件五级评分锚点：1 very slightly or not at all；2 a little；3 moderately；4 quite a bit；5 extremely。手册印刷页3的此刻时间窗口为right now；两组前后测使用同一固定时间窗口、题序和评分锚点。中文题文未核准前不把评分键当作可直接向校园学生施测的正式中文问卷。历史2003/2016译本与既有索取函保留为中文核对线索，不再承担英文原件许可的前置阻断。

\reviewsubsection{7.2 两条件中性教学脚本}
0-\allowbreak{}30秒：这是短时呼吸练习，舒适和自主停止优先；不要求完成所有步骤。

30-\allowbreak{}60秒：说明共同结构与当前/\allowbreak{}下一提示；漏步骤不追补，停留不用强求。

120-\allowbreak{}150秒：演示信号/\allowbreak{}节奏异常与自然呼吸收尾；暂停恢复从完整周期开始。

150-\allowbreak{}180秒：说明停止/\allowbreak{}求助方式，允许中性提问；不把答题正确作为入组条件。

A 60—120秒：圈是预设目标提示，场景按时间变化；它们不表示你的实际呼吸或达标，圈不会自动等你。舒适跟随或自然呼吸即可。

B 60—120秒：原生目标运动给出当前引导，独立实际运动只在信号可用时响应你的呼吸；信号不可用不等于失败。系统可有限等待/\allowbreak{}放慢并尝试回归，不能保证达标。累计是可靠练习响应，不是情绪恢复。

两条件60至120秒内容准确但不同；语速/\allowbreak{}示例数/\allowbreak{}操作者关注按同一预算。示意为设计教学，不声称正式180秒已可用。

\reviewsubsection{7.3 理解题与答案蓝图}
COMMON4题，A另3题，B另4题；共同题等值未验证，不能将各条件总分跨组排名。题材在情绪后测之后呈现，分别记录回应、跳过、超时、技术未呈现；后3类correct为空，不当错误。所有题仍为未验证设计题。

\reviewsubsection{7.4 C01 COMMON}
题干：画面明确给出当前吸气目标提示。这个提示能否单独证明你此刻确实在吸气？

构念：当前目标；答案要点：当前目标是吸气；不能仅凭目标提示确认自己的真实动作；真值：G的已呈现步骤。

计分：核心含义正确记1；错误含义记0；双人按答案要点编码，分歧第三人裁决；验证：两组相同文字任务，用各自正常载体；同构念不等于题项已经等值。

\reviewsubsection{7.5 C02 COMMON}
题干：引导进入标明自然呼吸的收尾。这意味着需要补做未完成周期吗？

构念：预算与停止；答案要点：不需要；自然呼吸收尾，可以停止；真值：R02和共同预算。

计分：核心含义正确记1；错误含义记0；双人按答案要点编码，分歧第三人裁决；验证：双方都使用相同收尾说明；B不额外延时。

\reviewsubsection{7.6 C03 COMMON}
题干：你漏过一个提示，随后已经开始呼气。现在应追补吸气吗？

构念：自主与漏步骤；答案要点：不用追补；舒适自然呼吸，按后续完整引导重新开始或停止；真值：R01/\allowbreak{}R06及Fade规则。

计分：核心含义正确记1；错误含义记0；双人按答案要点编码，分歧第三人裁决；验证：两组均讲解不追补，不测自己没有的信息。

\reviewsubsection{7.7 C04 COMMON}
题干：练习不舒服或想暂停，应该怎么做？

构念：自主停止；答案要点：停止/\allowbreak{}请求暂停，自然呼吸，不必等当前周期结束；真值：R01/\allowbreak{}R03。

计分：核心含义正确记1；错误含义记0；双人按答案要点编码，分歧第三人裁决；验证：共同停止程序与教学一致。

\reviewsubsection{7.8 A01 A}
题干：场景逐渐变清晰，是否说明你自己的呼吸已被确认达标？

构念：预设误读；答案要点：否；A场景按预定时间变化，不代表个人表现；真值：A输出独立X。

计分：核心含义正确记1；错误含义记0；双人按答案要点编码，分歧第三人裁决；验证：只在所属条件解释，不与另一条件专属题排名。

\reviewsubsection{7.9 A02 A}
题干：如果你的呼吸比圈慢，圈会自动等你吗？

构念：固定引导含义；答案要点：不会自动等；可以舒适跟随、自然呼吸或停止；真值：A开环。

计分：核心含义正确记1；错误含义记0；双人按答案要点编码，分歧第三人裁决；验证：只在所属条件解释，不与另一条件专属题排名。

\reviewsubsection{7.10 A03 A}
题干：你看到圈正常运动，能否从中判断传感器当前正常？

构念：设备与引导区分；答案要点：不能；A圈不表示设备状态；真值：A没有Actual/\allowbreak{}质量层。

计分：核心含义正确记1；错误含义记0；双人按答案要点编码，分歧第三人裁决；验证：只在所属条件解释，不与另一条件专属题排名。

\reviewsubsection{7.11 B01 B}
题干：目标风纹与实际风纹方向/\allowbreak{}进度不同，能否说你已经达到共同理想？

构念：目标与实际分离；答案要点：不能；目标G与真实X不同，是否符合I需独立评价；真值：I/\allowbreak{}G/\allowbreak{}X独立。

计分：核心含义正确记1；错误含义记0；双人按答案要点编码，分歧第三人裁决；验证：只在所属条件解释，不与另一条件专属题排名。

\reviewsubsection{7.12 B02 B}
题干：出现信号暂不可用，且实际风纹撤下。应该理解为呼吸失败吗？

构念：未知；答案要点：不能；当前不能可靠判断实际表现，可自然呼吸或停止；真值：R04/\allowbreak{}U未知。

计分：核心含义正确记1；错误含义记0；双人按答案要点编码，分歧第三人裁决；验证：只在所属条件解释，不与另一条件专属题排名。

\reviewsubsection{7.13 B03 B}
题干：引导放慢后你更容易跟上，是否已经证明更接近共同理想？

构念：支持与目标代价；答案要点：没有；更容易跟G不证明更接近I；真值：R07与独立I评分。

计分：核心含义正确记1；错误含义记0；双人按答案要点编码，分歧第三人裁决；验证：只在所属条件解释，不与另一条件专属题排名。

\reviewsubsection{7.14 B04 B}
题干：断流恢复后累计没有补涨，是否应该补做丢失步骤？

构念：累计/\allowbreak{}恢复；答案要点：不应补做；恢复不回填未知和未发生事件；真值：R04/\allowbreak{}R05及累计规则。

计分：核心含义正确记1；错误含义记0；双人按答案要点编码，分歧第三人裁决；验证：只在所属条件解释，不与另一条件专属题排名。

\newpage\reviewsection{八、访谈、体验表与地板决策}
\reviewsubsection{8.1 需求访谈与形成性记录}
1．最近一次学习间隙，你为什么想休息？实际用了什么办法？不要先介绍SRP优点。

2．是否遇到固定呼吸提示难跟随？请描述具体一次；没有经历也如实记录。

3．你希望短时练习帮助完成什么任务？怎样才算有用？

4．先展示A中性示意：你认为圈与场景变化分别代表什么？

5．再按所属条件展示正常/\allowbreak{}异常示意：哪里需要猜测、追赶或查看？

6．是否愿意佩戴两设备？准备、清洁、校准与操作员帮助有哪些顾虑？

7．如果自然休息、音频引导也能满足需要，你会选择哪一种，为什么？

8．什么体验会让你暂停或拒绝？再次自愿使用需要哪些条件？

8名需求访谈计入12名认知学生子集，不另加人数。先问已有经历再展示所属条件示意；记录无需求、拒绝及反例，不编码为项目失败者。体验后访谈追问具体困惑/\allowbreak{}追赶/\allowbreak{}支持/\allowbreak{}退出事件，不以偏爱解释情绪机制。

\reviewsubsection{8.2 自编体验与成本表}
EFFORT：刚才跟随提示需要你投入多少努力？；0 完全不费力／10 非常费力；用途：执行努力；缺失：未回答保留null。自编单题，未验证，不合成总分。

COMFORT：刚才的呼吸练习对你有多舒适？；0 完全不舒适／10 非常舒适；用途：呼吸舒适；缺失：未回答保留null。自编单题，未验证，不合成总分。

REUSE：如果下次学习间隙可自主选择，你愿意再次使用这套方案吗？请说明原因。；愿意／不愿意／不确定；用途：自愿使用意向，不代替真实复用；缺失：未回答保留null。自编产品问题，不合成总分。

成本记录分别填写：准备、佩戴、校准、教学、体验、后测/\allowbreak{}访谈、重置/\allowbreak{}清洁、暂停、故障/\allowbreak{}帮助、拒绝/\allowbreak{}退出的分钟/\allowbreak{}人时、帮助次数、退出时间与理由。愿意再次使用不是实际重复使用。

\reviewsubsection{8.3 低负性情绪的事前程序}
分别报前/\allowbreak{}后最低分比例及Wilson区间、可下降空间分布、缺题/\allowbreak{}耗时/\allowbreak{}理解；盲于组间效益。先认知审查指令；技术预试用整体地板/\allowbreak{}残差不确定域模拟主要差及精度。若目标关注差在可下降空间/\allowbreak{}资源下不可辨，不把群体筛高，收缩为估计研究或正式前重新论证量表；不自动替换终点。无普适地板百分数；以具体版本、关注差与可辨精度的先验决策为依据，当前无真人数据。

\newpage\reviewsection{九、完整分析规格与正式N确定程序}
\reviewsubsection{9.1 估计对象与主要模型}
同一规定时间预算下B减A的NA分配效应；原分配全部入目标人群。

NA\_\allowbreak{}post\textasciitilde{}B+centered\_\allowbreak{}NA\_\allowbreak{}pre+四顺序固定效应；HC3，双侧95\%区间。

\reviewsubsection{9.2 缺失与模型合并}
候选MICE预测均值匹配k=5，50插补、每条20迭代；仅设计工作参数须用合成及前置数据验证。条件/\allowbreak{}基线/\allowbreak{}顺序/\allowbreak{}预记录经验/\allowbreak{}记录时延作为候选预测；基线缺失也联合插补，不以均值填充。

每插补同一模型与HC3，再Rubin合并；固定seed=20261008用于设计复算，不是运行随机化seed。

\reviewsubsection{9.3 缺失敏感性}
完全案例；按组对已插补NA施加delta，范围[-\allowbreak{}MID,+MID]及±2MID、截至确切版本允许范围，报告翻转区域；MID未知时不伪填数。

\reviewsubsection{9.4 次要与多重性}
独立I连续偏差与C/\allowbreak{}K/\allowbreak{}U界、G误差、控制成本、努力/\allowbreak{}舒适。估计及95\%区间，不用次要p值补主要失败；若确认性次要检验另冻结Holm族。

\reviewsubsection{9.5 正式规模确定}
MID与残差/\allowbreak{}相关/\allowbreak{}地板/\allowbreak{}缺失区间组合，模拟同一模型和缺失方案；每场景2000次为计算预算候选，报告MonteCarlo误差；给80/\allowbreak{}90\%功效及区间半宽，不倒推MID；没有真实输入不产正式N。

NA前测缺失不以事后排样规避；缺失方法与参数需前置验证再正式冻结。原始X不由G插补。顺序为四天气固定效应；ANCOVA没有处理后美感/\allowbreak{}努力/\allowbreak{}适应次数协变量。主模型无法估计时报告不足，不自动换终点或删顺序；预先模拟须检验该情景。

\reviewsubsection{9.6 校准、冻结及结论的统一顺序}
设计定稿→对应前置活动批准与材料固定→前置实现及验证→正式参数、分析与协议冻结→正式研究。前置活动有自己的材料/\allowbreak{}程序固定，不等于正式参数先于预试冻结。

两种100人资源情景保留；20技术前置按10校准/\allowbreak{}10保留核验，若新增前置人员正式可用人数减少。四顺序8/\allowbreak{}16隐藏块仅下一版候选，运行旧算法不改。M01/\allowbreak{}M02未关闭时不能形成可施测主要结果材料。正式协议、真实批准、数值与N均未由本设计取得。



\newpage\reviewsection{十、中文工作题本与核准程序}
\reviewsubsection{10.1 中文审阅工作稿}
版本SRP-PANASX-GENERAL-ZH-DRAFT-01（2026年10月9日）。本稿由AI直接依据PANAS-X英文一般分量表起草，不是官方中文版本，不冒称黄丽等2003或郭明珠、甘怡群2010译本；尚无独立语言审阅、回译或真人验证。英文原件的非商业复制使用已有依据，中文翻译/改编及公开范围仍需针对实际用途核对；当前用于内部审阅。本轮没有施测。

指导语工作稿：下面列出一些描述感受的词。请根据你此刻的实际感受，选择每种感受的强烈程度。没有正确或错误答案；请不要按平时的状态或你认为练习应该带来的变化作答。若不愿回答，可以留空。

评分锚点工作稿：1 几乎没有或完全没有；2 有一点；3 中等程度；4 相当强烈；5 非常强烈。每题单选1至5，未回答保留缺失，不强制作答。以下英文及分量归属供审阅者核对，正式参与者表只展示核准中文与选项；题序为项目工作顺序，两组前后测一致。

\begin{longtable}{L{12mm}L{33mm}L{36mm}L{16mm}L{40mm}}
编号 & 英文原项 & 中文工作措辞 & 归属 & 记录\\\hline
Q01 & interested & 感兴趣的 & PA & 1 / 2 / 3 / 4 / 5 / 留空\\
Q02 & distressed & 苦恼的 & NA & 1 / 2 / 3 / 4 / 5 / 留空\\
Q03 & excited & 兴奋的 & PA & 1 / 2 / 3 / 4 / 5 / 留空\\
Q04 & upset & 心烦的 & NA & 1 / 2 / 3 / 4 / 5 / 留空\\
Q05 & strong & 强有力的 & PA & 1 / 2 / 3 / 4 / 5 / 留空\\
Q06 & guilty & 内疚的 & NA & 1 / 2 / 3 / 4 / 5 / 留空\\
Q07 & scared & 惊恐的 & NA & 1 / 2 / 3 / 4 / 5 / 留空\\
Q08 & hostile & 有敌意的 & NA & 1 / 2 / 3 / 4 / 5 / 留空\\
Q09 & enthusiastic & 热情的 & PA & 1 / 2 / 3 / 4 / 5 / 留空\\
Q10 & proud & 自豪的 & PA & 1 / 2 / 3 / 4 / 5 / 留空\\
Q11 & irritable & 易怒的 & NA & 1 / 2 / 3 / 4 / 5 / 留空\\
Q12 & alert & 警觉的 & PA & 1 / 2 / 3 / 4 / 5 / 留空\\
Q13 & ashamed & 羞愧的 & NA & 1 / 2 / 3 / 4 / 5 / 留空\\
Q14 & inspired & 受到鼓舞的 & PA & 1 / 2 / 3 / 4 / 5 / 留空\\
Q15 & nervous & 紧张的 & NA & 1 / 2 / 3 / 4 / 5 / 留空\\
Q16 & determined & 坚定的 & PA & 1 / 2 / 3 / 4 / 5 / 留空\\
Q17 & attentive & 专注的 & PA & 1 / 2 / 3 / 4 / 5 / 留空\\
Q18 & jittery & 忐忑不安的 & NA & 1 / 2 / 3 / 4 / 5 / 留空\\
Q19 & active & 活跃的 & PA & 1 / 2 / 3 / 4 / 5 / 留空\\
Q20 & afraid & 害怕的 & NA & 1 / 2 / 3 / 4 / 5 / 留空\\
\end{longtable}
NA：Q02、Q04、Q06、Q07、Q08、Q11、Q13、Q15、Q18、Q20；PA为其余10题。每个分量独立求和10至50，无反向题，不计算20题总分。任一分量缺题则该分量总分缺失，另报缺题数；这是项目工作缺题规则。合成检查：全1则各10，全5则各50；只缺Q02则NA缺失、PA不受影响。检查只证明映射与求和规则，不证明测量有效。

\reviewsubsection{10.2 必须核对的语义与独立审阅}
重点核查：ashamed指羞愧，不以社交害羞替代；strong是否被只理解为体力，alert是否被只理解为危险戒备；jittery与nervous是否都仅被回答为紧张；upset与distressed是否被误认同一词；scared与afraid是否引入原版没有规定的固定强弱排序；inspired是否被理解为创造灵感而非受鼓舞；active是否被理解为活动次数而非此刻感受。上述是待检验疑点，不是已发生的用户误读，不在正式答题时追加诱导解释。

F-02与U12-04负责逐题语言和构念审阅，Q-01承接合适评审者，Q-02承接认知材料。先由具备中英语言与测量知识者独立评阅题项、锚点和指导语；另一人仅看中文进行回译，再对照原项讨论。AI起草不能计作两名独立译者。优先在原8位评审内安排2位合适人员；若没有相应能力，另列人员成本，不假定已满足。

逐题填写：编号、受访者复述、具体例子、时间窗理解、选项使用理由、误读类型、拟修订词、复核结论。只复述指导语，不向被试教标准情绪答案。请其区分“此刻”与“最近一周”，说明1与5的含义，并解释实际选择；不要求其报告不愿公开的个人经历。

12名认知学生内部安排前6人发现问题、后6人核查修订，每轮兼顾两条件，仍计入现有12人，不额外加人数。每个改变构念或时间窗的误读必须有处理和后续独立核查；重大语义分歧不能靠多数票直接消除。若最后一轮仍需实质改词，增加核查资源并重算正式预算，不用同一批已被解释答案者冒充独立核验。

\reviewsubsection{10.3 通过、失败与研究边界}
材料准入条件：20题、指导语和五级锚点均完成审阅；重大语义分歧无未决项；最后一轮没有未处理的构念/时间窗误读；F-02/U12-04记录采用理由，G-05核对实际翻译及使用范围。这是项目材料核准规则，不是通用效度阈值。方法参考\href{https://www.intestcom.org/files/guideline_test_adaptation_2ed.pdf}{ITC 2017测试翻译与适配指南}的语言复核及经验验证要求；具体人员与轮次为本项目资源方案。

20名技术前置用于耗时、缺失、地板及分布检查，不能仅凭小样本内部一致性或回译通过声称新译本已验证。正式前由F-02/U12-04根据选定版本及已有证据确定是否需要额外独立测量验证样本；若需要，单独预算，不占用同一正式样本作选词和确认效应。不能按A/B情绪差异选词、删除题项或选择量表版本。

中文材料或测量证据不足时暂停对应正式施测；已有合适中文版本若随后取得，逐项比较后只选一个版本，不拼接。正式固定后，两组前后测使用完全一致的版本；修改后另起版本并重审，不能合并不同版本充当前后变化。selected\_form与正式item\_texts、设计全面闭合状态均不因本稿完成而自动填写。论文可称“依据PANAS-X一般分量表形成的中文适配稿”，未取得证据前不能称“经验证中文版PANAS”。


\newpage\reviewsection{十一、前置审阅执行表}
\reviewsubsection{11.1 交接对象与输入}
本节为设计批准后的执行准备材料，空表不是执行结果。F-02统筹题本，U12-04核对测量与评分，Q-01安排语言审阅，Q-02安排认知访谈，G-05核对相应活动和材料使用范围。实际人员、活动日期及活动批准尚待填写。输入为第十节工作稿、PANAS-X原手册印刷页1--3及本节记录表；输出回填本稿/设计索引，不建立另一套规范。

回译者仅接收第十节中文题项、中文指导语和中文选项，不接收英文、NA/PA归属或预期译法；由协调人提供中文摘录。完成回译后再打开英文原件对照。协调人记录是否事先见过原项；见过原项的审阅仍可参考，但不能标为盲于原文的独立回译。

\reviewsubsection{11.2 语言审阅登记及逐项复核表}
材料版本：\rule{62mm}{0.4pt}\quad 日期：\rule{30mm}{0.4pt}

正向审阅人/语言与测量经验：\rule{88mm}{0.4pt}

回译人/语言经验：\rule{104mm}{0.4pt}

回译前是否见过英文原项：是／否／不确定；协调人：\rule{40mm}{0.4pt}

每行填“保留／修改／未决”，理由必须指向词义或构念，不按字面相似度评分。下表可逐行扩展；英文回译与讨论须先分别完成，再合议。说明与锚点同样审阅，不只检查20个词。
\begin{longtable}{L{16mm}L{29mm}L{29mm}L{42mm}L{26mm}}
对象 & 独立回译 & 分歧/可能误读 & 处理理由与拟定措辞 & 结论/复核人\\\hline
Q01 & & & & \\[5mm]
Q02 & & & & \\[5mm]
Q03 & & & & \\[5mm]
Q04 & & & & \\[5mm]
Q05 & & & & \\[5mm]
Q06 & & & & \\[5mm]
Q07 & & & & \\[5mm]
Q08 & & & & \\[5mm]
Q09 & & & & \\[5mm]
Q10 & & & & \\[5mm]
Q11 & & & & \\[5mm]
Q12 & & & & \\[5mm]
Q13 & & & & \\[5mm]
Q14 & & & & \\[5mm]
Q15 & & & & \\[5mm]
Q16 & & & & \\[5mm]
Q17 & & & & \\[5mm]
Q18 & & & & \\[5mm]
Q19 & & & & \\[5mm]
Q20 & & & & \\[5mm]
指导语 & & & & \\[5mm]
五级锚点 & & & & \\[5mm]
\end{longtable}
\reviewsubsection{11.3 认知访谈主持稿及人员安排}
启动前核对对应活动批准、知情材料、使用范围和已完成语言审阅的题本版本；缺项填写“待补”，不把设计批准当作真人活动批准。第一轮6名、第二轮6名不同学生，每轮A/B各3名是材料覆盖安排，不是效益比较随机化；累计仍为A/B各6名。需求访谈的8人仍是这12人子集。实际到场、退出与替补分别登记，不把邀请数当完成人数。

主持开场：“我们想检查这些词和说明是否清楚，不是在考你。请按自己的理解回答；不明白可以直接说，也可以跳过或停止。”

步骤1：在观看条件教学前，让参与者自行阅读中文工作题本并回答；不展示英文或正负情绪标签。记录开始/结束时间、跳过和求助，不设置未经试验的限时。

步骤2：先问“这份说明让你回想哪个时间段？”“1和5分别是什么意思？”再逐题问“你怎样理解这个词？可以用一句话或一个虚构例子说明。”“你选择这个选项时考虑了什么？”出现含混时只问“可以再具体一点吗？”，不先给标准解释。

步骤3：对相近词问“这两个词对你是一样还是不同？请说明。”不要求一定存在差别或强弱次序。对不愿分享的感受，不追问个人经历。正常情绪低分不算理解失败。

步骤4：按所属条件进行既定教学/故事板理解材料。第一轮结束后汇总再改词，第二轮只用已固定的修订版本；不能对每个参与者临时换词。认知访谈包含反复追问，因此其情绪答案不进入正式效益分析。

\reviewsubsection{11.4 单人认知记录表（每人一份）}
匿名编号：\rule{28mm}{0.4pt}\quad 轮次：1／2\quad 条件：A／B

题本版本：\rule{43mm}{0.4pt}\quad 主持人/日期：\rule{48mm}{0.4pt}

对应活动批准引用：\rule{97mm}{0.4pt}

开始/结束时间：\rule{43mm}{0.4pt}\quad 完成／退出／未开始

逐题复制下列记录格，覆盖Q01--Q20、指导语及锚点。空白代表未记录，不能记为“没有误读”。不在表内登记姓名、联系方式；身份管理按相应活动材料执行。
\begin{longtable}{L{39mm}L{110mm}}
字段 & 待填写内容\\\hline
对象编号/原始选择 & \rule{85mm}{0.4pt}\\[5mm]
原话复述/例子 & \rule{85mm}{0.4pt}\\[9mm]
选择理由/时间窗 & \rule{85mm}{0.4pt}\\[9mm]
求助/跳过及说明 & \rule{85mm}{0.4pt}\\[5mm]
理解分类与依据 & 保留原意／措辞不自然／构念偏移／时间窗偏移／锚点误读／未能判断；依据：\\[9mm]
主持人是否提示解释 & 无／有（记录原话，解释后回答不算独立核验）\\[5mm]
修订建议/复核人 & \rule{85mm}{0.4pt}\\[5mm]
\end{longtable}
\reviewsubsection{11.5 轮次汇总与决定}
F-02汇总每个问题涉及的人次、原话、严重性与处理；分母为实际接受该题核查的人数，未核查者不计为通过。U12-04复核构念、计分与证据边界。保留所有相反例，不用平均满意度覆盖重大误读。

决定表：轮次/材料版本：\rule{45mm}{0.4pt}；实际完成数A：\rule{10mm}{0.4pt}、B：\rule{10mm}{0.4pt}；未决问题编号：\rule{50mm}{0.4pt}；拟修订/理由：\rule{60mm}{0.4pt}；新版本/下一轮人员需求：\rule{45mm}{0.4pt}；复核人/日期：\rule{45mm}{0.4pt}。

允许进入下一步：逐项语言审阅和使用范围核对完成、重大分歧有处理、最后一轮无未处理构念或时间窗误读，才可固定前置流程测量材料。退回修订：任一重大分歧未决、未经独立核查的实质改词或记录不足；补充材料/人员后再核查。正式主要测量的选定仍需适用性证据评估，不由这张表自动放行。

交付验收分为三栏：材料准备（本节已完成）；语言与认知执行（待执行）；正式测量选定（待论证与冻结）。下一执行入口为F-02/Q-01的实际审阅人员落实和中文摘录准备，而不是招募正式研究参与者。任何实际签名、回译、访谈回答均由执行后填写，本稿不预填通过结果。

\end{document}
